% command to make sure colorboxes don't create whitespace between code lines
\fboxsep=.6pt

% default settings for listings
\lstset{style=python, basicstyle=\ttfamily\footnotesize}

\title{Theoretische Informatik}

\subtitle{Kolmogorov-Komplexität}

\author[Graf, Thomas] % (optional)
{Thomas Graf}

\institute[KLW] % (optional)
{
	Fachschaft Informatik\\
	Kantonsschule im Lee
}

\date[\today] % (optional)

\logo{\includegraphics[width=1cm]{Logos/FreeFlower.pdf}}

\titlegraphic{\vspace{-.7cm}\includegraphics[width=\textwidth]{"Grundlagen_Info/06_Datenbanken/Figures/title-emojis"}}

\frame{\titlepage}

\begin{frame}[fragile]{Definition der Kolmogorov-Komplexität}
\begin{mydefinition}[Kolmogorov-Komplexität]\label{mydefinition:kolmogorov}
Es sei $w$ ein binäres Wort. Die plain Kolmogorov-Komplexität $C(w)$ von $w$ ist definiert als die \textbf{Anzahl der Bits eines kürzesten Maschinencodes}, der durch Kompilierung eines \texttt{C'++}-Programms mit \texttt{g'++} entsteht und dessen Programm auf der Eingabe $\lambda$ (dem leeren Wort) genau $w$ ausgibt und anschliessend hält. Wir sagen dann, dass das Programm $w$ generiert.
\end{mydefinition}
Für ein binäres Wort $w$ betrachten wir alle mit \texttt{g'++} erzeugten Maschinencodes von \texttt{C'++}-Programmen, die $w$ generieren. Die Länge eines kürzesten\footnote{Im Allgemeinen kann es mehrere verschiedene kürzeste Maschinencodes geben, deren Programme $w$ generieren.} solchen Maschinencodes ist dann die Zahl $C(w)$.
\end{frame}

\begin{frame}{Klassen in Python}
ok
\end{frame}

\begin{frame}[fragile]{Listen als Klassen}
ok
\end{frame}

\setbeamertemplate{logo}{} % Hide logo


\begin{frame}[fragile]{Beispiel: Personen}
ok
\end{frame}
